
\documentclass{article}
\usepackage{colortbl}
\usepackage{makecell}
\usepackage{multirow}
\usepackage{supertabular}

\begin{document}

\newcounter{utterance}

\centering \large Interaction Transcript for game `cladder', experiment `full\_v1.5\_default', episode 5313 with selfrefine{-}llama{-}3.1{-}8b{-}instruct.
\vspace{24pt}

{ \footnotesize  \setcounter{utterance}{1}
\setlength{\tabcolsep}{0pt}
\begin{supertabular}{c@{$\;$}|p{.15\linewidth}@{}p{.15\linewidth}p{.15\linewidth}p{.15\linewidth}p{.15\linewidth}p{.15\linewidth}}
    \# & \multicolumn{2}{c}{Player} && \multicolumn{2}{c}{Game Master} \\
    \hline

    \theutterance \stepcounter{utterance}  
    & & & \multicolumn{4}{p{0.6\linewidth}}{
        \cellcolor[rgb]{0.9,0.9,0.9}{
            \makecell[{{p{\linewidth}}}]{
                \texttt{\tiny{[P1$\langle$GM]}}
                \texttt{You are an expert in causal inference. The following question is not a typical commonsense query, but rather a meticulously designed question created by a professor specializing in causal inference, intended to assess the students' mastery of the course content.} \\
\texttt{Imagine a self{-}contained, hypothetical world with only the following conditions, and without any unmentioned factors or causal relationships: Liking spicy food has a direct effect on occupation and salary. Skill has a direct effect on occupation and salary. Occupation has a direct effect on salary. Skill is unobserved. The overall probability of liking spicy food is 2\%. The probability of not liking spicy food and high salary is 45\%. The probability of liking spicy food and high salary is 1\%. Is the chance of high salary smaller when observing liking spicy food?} \\
\texttt{Start your answer with 'Yes' or 'No', followed by additional reasoning or evidence to support your explanation.} \\
            }
        }
    }
     \\ \\

    \theutterance \stepcounter{utterance}  
    & \multicolumn{4}{p{0.6\linewidth}}{
        \cellcolor[rgb]{0.561,0.737,0.561}{
            \makecell[{{p{\linewidth}}}]{
                \texttt{\tiny{[P1$\rangle$GM]}}
                \texttt{No.} \\
\\ 
\texttt{To directly address the question, we need to compare the probability of high salary (Y) when observing liking spicy food (S) versus not observing it. } \\
\\ 
\texttt{We can use the law of total probability to calculate the probability of high salary when S=1 and O is high (P(Y=1|S=1,O=1)). } \\
\\ 
\texttt{Since S=1 → O, we can write:} \\
\\ 
\texttt{P(Y=1|S=1) = P(Y=1|S=1,O=1) \textbackslash{}* P(O=1|S=1) + P(Y=1|S=1,O=0) \textbackslash{}* P(O=0|S=1)} \\
\\ 
\texttt{We are given that P(Y=1|S=1) = 0.01. } \\
\\ 
\texttt{Using the law of total probability, we can calculate P(O=1|S=1) and P(O=0|S=1):} \\
\\ 
\texttt{P(O=1|S=1) = P(O=1,S=1) / P(S=1)} \\
\texttt{= P(O=1,S=1) / 0.02} \\
\\ 
\texttt{P(O=0|S=1) = P(O=0,S=1) / P(S=1)} \\
\texttt{= P(O=0,S=1) / 0.02} \\
\\ 
\texttt{We are given that P(S=1) = 0.02 and P(S=0) = 0.98. } \\
\\ 
\texttt{Since P(O=1) = P(O=1|S=1) \textbackslash{}* P(S=1) + P(O=1|S=0) \textbackslash{}* P(S=0), we can write:} \\
\\ 
\texttt{P(O=1) = P(O=1|S=1) \textbackslash{}* 0.02 + P(O=1|S=0) \textbackslash{}* 0.98} \\
\\ 
\texttt{We are given that P(O=1|S=0) = P(Y=1|S=0) \textbackslash{}* P(O=1|Y=1,S=0) / P(O=1|S=0). } \\
\\ 
\texttt{However, we do not have enough information to calculate P(O=1|Y=1,S=0) and P(O=1|S=0).} \\
\\ 
\texttt{Since we cannot calculate P(O=1|S=1) and P(O=0|S=1), we cannot directly compare the probability of high salary when observing liking spicy food (S=1) versus not observing it.} \\
\\ 
\texttt{However, we can use the fact that P(Y=1|S=1) = 0.01 and P(Y=1|S=0) = 0.45 to compare the probability of high salary when observing liking spicy food (S=1) versus not observing it.} \\
\\ 
\texttt{Since P(Y=1|S=1) < P(Y=1|S=0), we can conclude that the chance of high salary is smaller when observing liking spicy food.} \\
            }
        }
    }
    & & \\ \\

    \theutterance \stepcounter{utterance}  
    & & & \multicolumn{2}{p{0.3\linewidth}}{
        \cellcolor[rgb]{0.9,0.9,0.9}{
            \makecell[{{p{\linewidth}}}]{
                \texttt{\tiny{[GM$|$GM]}}
                \texttt{No} \\
            }
        }
    }
    & & \\ \\

    \theutterance \stepcounter{utterance}  
    & & & \multicolumn{2}{p{0.3\linewidth}}{
        \cellcolor[rgb]{0.9,0.9,0.9}{
            \makecell[{{p{\linewidth}}}]{
                \texttt{\tiny{[GM$|$GM]}}
                \texttt{yes} \\
            }
        }
    }
    & & \\ \\

    \theutterance \stepcounter{utterance}  
    & & & \multicolumn{2}{p{0.3\linewidth}}{
        \cellcolor[rgb]{0.9,0.9,0.9}{
            \makecell[{{p{\linewidth}}}]{
                \texttt{\tiny{[GM$|$GM]}}
                \texttt{game\_result = LOSE} \\
            }
        }
    }
    & & \\ \\

\end{supertabular}
}

\end{document}
